\documentclass[dvisvgm,tikz,border=3pt]{standalone}
\usepackage{amsmath,amssymb}
\usepackage{pgfplots}
\usepackage{CJKutf8}
\pgfplotsset{compat=1.18}
\usetikzlibrary{arrows.meta,calc,positioning}

\definecolor{themebg}{HTML}{30362d}
\definecolor{themefg}{HTML}{edf4e8}
\definecolor{themegray}{HTML}{9ea897}
\definecolor{themecurve}{HTML}{7eb6ff}
\definecolor{themealert}{HTML}{ff7b7b}
\definecolor{themeamber}{HTML}{f5b942}

\pgfdeclarelayer{background}
\pgfdeclarelayer{foreground}
\pgfsetlayers{background,main,foreground}

\begin{document}
\begin{CJK*}{UTF8}{gbsn}
\pagecolor{themebg}
\begin{tikzpicture}[color=themefg, text=themefg, >=Stealth]

% ── 上半部：映射级联流水线 (加宽间距，彻底消除文字与框线重合) ──
\begin{scope}[shift={(0,0)}]
  % 节点 1：输入空间
  \node[draw=themegray!70, fill=themebg, rounded corners=5pt, line width=1.0pt, inner sep=6pt, font=\bfseries\small] (In1) at (0, 0) {输入空间 $\mathbb{R}^n$};

  % 节点 2：右乘 Q 输入重组
  \node[draw=themeamber, fill=themeamber!15!themebg, text=themeamber, rounded corners=5pt, line width=1.2pt, inner sep=6pt, font=\bfseries\small] (In2) at (4.4, 0) {输入重组 $\mathbb{R}^n$};

  % 节点 3：核心映射 A
  \node[draw=themecurve, fill=themecurve!15!themebg, text=themecurve, rounded corners=5pt, line width=1.2pt, inner sep=6pt, font=\bfseries\small] (Out1) at (8.8, 0) {核心输出 $\mathbb{R}^m$};

  % 节点 4：左乘 P 输出重组
  \node[draw=themealert, fill=themealert!15!themebg, text=themealert, rounded corners=5pt, line width=1.2pt, inner sep=6pt, font=\bfseries\small] (Out2) at (13.2, 0) {输出重组 $\mathbb{R}^m$};

  % 连接有向边与清晰注记 (带背景遮罩，防止穿线)
  \draw[->, themeamber, line width=2.0pt] (In1) -- 
    node[above=3pt, font=\bfseries\small, text=themeamber, fill=themebg, inner sep=1.5pt] {右乘 $\boldsymbol{Q}$} 
    node[below=3pt, font=\scriptsize, text=themegray, fill=themebg, inner sep=1.5pt] {可逆列变换} (In2);
  
  \draw[->, themecurve, line width=2.4pt] (In2) -- 
    node[above=3pt, font=\bfseries\small, text=themecurve, fill=themebg, inner sep=1.5pt] {映射 $\boldsymbol{A}$} 
    node[below=3pt, font=\scriptsize, text=themegray, fill=themebg, inner sep=1.5pt] {尺寸 $m \times n$} (Out1);

  \draw[->, themealert, line width=2.0pt] (Out1) -- 
    node[above=3pt, font=\bfseries\small, text=themealert, fill=themebg, inner sep=1.5pt] {左乘 $\boldsymbol{P}$} 
    node[below=3pt, font=\scriptsize, text=themegray, fill=themebg, inner sep=1.5pt] {可逆行变换} (Out2);

  % 顶层跨越弧线：总复合 PAQ
  \draw[->, dashed, themefg, line width=1.2pt] (In1.north) to[bend left=20] 
    node[above=2pt, font=\bfseries\small, fill=themebg, inner sep=2pt] {总复合映射：$\boldsymbol{P}\boldsymbol{A}\boldsymbol{Q}$（保持 $\operatorname{rank}$ 等价）} (Out2.north);
\end{scope}

% ── 下半部：左右乘机理对比卡片 ──
\begin{scope}[shift={(0,-2.4)}]
  % 左侧卡片：左乘 P (行操作)
  \node[draw=themealert!60!themebg, fill=themebg, rounded corners=6pt, line width=0.8pt, inner sep=8pt, text width=6.2cm, anchor=north] (CardL) at (3.1, 0) {
    \centering
    {\bfseries\color{themealert} 左乘 $\boldsymbol{P}$：行操作 / 组装输出} \\[5pt]
    \raggedright\small
    $\bullet$ \textbf{行视角}：$(\boldsymbol{P}\boldsymbol{A})_i = \sum_k p_{ik}\boldsymbol{\alpha}_k$（第 $i$ 行是 $\boldsymbol{A}$ 各行的线性组合）\\
    $\bullet$ \textbf{几何含义}：对输出空间做基变换 / 坐标重组\\
    $\bullet$ \textbf{初等行变换}：消元解方程、化行最简阶梯形\\
    $\bullet$ \textbf{不变性}：保持行空间 $\operatorname{Row}(\boldsymbol{P}\boldsymbol{A}) = \operatorname{Row}(\boldsymbol{A})$（$\boldsymbol{P}$ 可逆）
  };

  % 右侧卡片：右乘 Q (列操作)
  \node[draw=themeamber!60!themebg, fill=themebg, rounded corners=6pt, line width=0.8pt, inner sep=8pt, text width=6.2cm, anchor=north] (CardR) at (10.1, 0) {
    \centering
    {\bfseries\color{themeamber} 右乘 $\boldsymbol{Q}$：列操作 / 重组输入} \\[5pt]
    \raggedright\small
    $\bullet$ \textbf{列视角}：$(\boldsymbol{A}\boldsymbol{Q})_j = \sum_k q_{kj}\boldsymbol{a}_k$（第 $j$ 列是 $\boldsymbol{A}$ 各列的线性组合）\\
    $\bullet$ \textbf{几何含义}：对输入空间做基变换 / 坐标重组\\
    $\bullet$ \textbf{初等列变换}：简化列向量生成组、化相抵标准形\\
    $\bullet$ \textbf{不变性}：保持列空间 $\operatorname{Col}(\boldsymbol{A}\boldsymbol{Q}) = \operatorname{Col}(\boldsymbol{A})$（$\boldsymbol{Q}$ 可逆）
  };
\end{scope}

% ── 底部总结横条 ──
\node[font=\bfseries\small, color=themecurve, anchor=north] at (6.6, -7.2) {
  核心心智模型：左乘重组输出/行，右乘重组输入/列；初等变换本质是在两端独立进行可逆基变换
};

\end{tikzpicture}
\end{CJK*}
\end{document}
